\documentclass[11pt]{article}
\usepackage[T1]{fontenc}
\usepackage{lmodern}
\usepackage[margin=1in]{geometry}
\usepackage{amsmath,amssymb,amsthm}
\numberwithin{equation}{section}
\usepackage[expansion=false]{microtype}
\usepackage[hidelinks,hypertexnames=false]{hyperref}
\setlength{\emergencystretch}{2em}
\newcommand{\workbenchurl}{https://github.com/KokunoYumeto/zeta-and-l-functions-workbench}
\newcommand{\recordurl}{https://github.com/KokunoYumeto/mathematics-workbenches}
\newcommand{\N}{\mathrm N}
\newcommand{\Dc}{\mathcal D}
\newcommand{\cont}{\operatorname{cont}}
\newtheorem{theorem}{Theorem}[section]
\newtheorem{lemma}[theorem]{Lemma}
\title{A smaller balanced prime-sum criterion\\for the cubic Riemann hypothesis\\
\large After Heath--Brown's 1982 identity}
\author{}
\date{24 November 2026}
\begin{document}
\maketitle
\begin{abstract}
Heath--Brown's 1982 identity expresses prime coefficients through
balanced inverse convolutions. We prove a rational-content estimate
for its cubic Hecke application that uses any specified zero-free
half-plane. Keeping ideal norm and rational content in one Euler
product yields a uniform tail bound, with the dependence on both
norm caps explicit. At square-root error this removes a nonempty
region from the balanced convolution. Applying Liu's stated
zero-free boundary to both the cubic Hecke function and zeta gives a
content exponent \(0.2849927501\ldots\), or
\(0.2494473999\ldots\) under classical
zeta RH, compared with the elementary values \(1/3\) and \(2/7\).
In the latter case the remaining primitive norm product has exponent
\(0.5011052001\ldots\), strictly greater than one half. We prove
that the improved cutoff removes nonzero summands present under the
older cutoff. Extracting one additional Hecke Euler factor also
proves the tail estimate throughout the full stated range of both
zero-free parameters. All finite factors, contour bounds and the exact
balanced identity are included. Liu's analytic zero-free theorem
is an external input; no further improvement of that zero-free
half-plane is claimed. The remaining all-weight square-root bound
is equivalent to cubic RH and is not proved. The mathematics and
text were produced by GPT-6 Astra in ChatGPT Pro mode, with separate
same-model re-proofs and targeted checks; no human has checked this work.
\end{abstract}
\section{Statements and scope}

The identity used below is the classical binomial identity of
Heath--Brown \cite[p.~1367, Lemma~1, equations~(6)--(7)]{JCTHeathBrown}.
The estimate concerns the part of the resulting balanced sum with
large product of rational contents. Keeping norm and content in one
Euler product converts any specified zero-free half-plane into a
smaller cutoff. In particular the published input of
Liu \cite[Theorems~1.2 and~8.29]{JCTLiu} improves the elementary
cutoffs below. The classical identity, the external zero-free theorem,
and the Mellin argument proved here have distinct roles.
The estimate removes a nonempty region at square-root cost; it does
not bound the signed sum that remains.

\paragraph{Fixed arithmetic and standard analytic background.}
Let \(K=\mathbf Q(\omega)\), where \(\omega^2+\omega+1=0\), and
let \(\psi\) be either faithful cubic character defining
\(K(\sqrt[3]{269})/K\), of conductor \((269)\).
Its conjugate under complex conjugation is \(\psi^{-1}\).
In particular \(\psi((r))=1\) for every rational integer
\(r\) prime to the conductor: that value equals its inverse and
has order dividing three.
Write \(L(s)=L_K(s,\psi)\).  Its rational conductor is
\(N=|D_K|\N(269)=3\cdot269^2=217083\).
The standard nontrivial finite-order Hecke theory supplies an
entire function of finite order, polynomial growth in every fixed
vertical strip, and the functional equation
\begin{equation}
 L(s)=\epsilon a_0^{2s-1}
       \frac{\Gamma(1-s)}{\Gamma(s)}L^\vee(1-s),
 \qquad a_0=\frac{2\pi}{\sqrt N},\qquad |\epsilon|=1.
 \label{jct:fe}
\end{equation}
Here \(L^\vee=L_K(s,\overline\psi)\); conjugating ideals gives
\(L^\vee=L\).  These are the general Hecke analytic background inputs; see
\cite[\S2.1]{JCTZaman}. The reciprocal bounds below are proved
from zero-freeness and strip growth, without a zero-count theorem.
Any stronger zero-free input is separately specified below.

Let \(S\) consist of the prime ideals above \(3\) and \(269\), and put
\begin{equation}
 P(s)=L_K^S(s,\psi)=H(s)L(s),\qquad
 H(s)=1-\psi(\lambda)3^{-s},\quad \lambda=\omega-1.
 \label{jct:incomplete}
\end{equation}
The prime above \(269\) already has no Euler factor in \(L\).
Unless explicitly stated otherwise, all ideals in sums below are
nonzero integral ideals prime to \(S\).
Set
\[
 m(\mathfrak a)=\mu_K(\mathfrak a)\psi(\mathfrak a),\qquad
 D(s)=P(s)^{-1}.
\]
For an ideal \(\mathfrak a\), its rational content
\(\cont(\mathfrak a)\) is the largest positive integer \(r\) such
that \(\mathfrak a=(r)\mathfrak b\) with \(\mathfrak b\) integral.
The latter ideal is primitive when its content is one.  Since
\(\mathbf Z[\omega]\) is a PID, a generator \(a+b\omega\) has
content \(\gcd(a,b)\), which is unchanged by multiplication by a unit.

Fix smooth functions \(b,\chi,\chi_d:[0,\infty)\to[0,1]\), each
equal to one on \([0,1]\) and zero on \([2,\infty)\).
Write \(\Dc=t\,d/dt\).  Constants may depend on these fixed
functions and on the fixed field and conductor.

\paragraph{The explicit zero-free input.}
Fix \(1/2<\beta\le1\) and \(1/2\le\rho\le1\), and assume
\begin{equation}
 P(s)\ne0\quad(\Re s>\beta),\qquad
 \zeta\text{ has no zero in }\Re s>\rho.
 \label{jct:zero-free-input}
\end{equation}
The values \(\beta=\rho=1\) follow from the absolutely convergent
Euler products and need no additional theorem. Classical zeta RH
permits \(\rho=1/2\). Put \(u=\rho/2\), and define
\begin{equation}
 a_\beta=2\beta-u,\qquad
 \gamma_{\beta,u}=\frac{\beta-1/2}{2\beta-u}.
 \label{jct:exponents}
\end{equation}
These parameters are fixed before the asymptotic variables.
There is no additional condition \(\beta+u\ge1\): the secondary
Euler factor extracted in the proof removes that restriction.

\begin{theorem}[Joint-content tail]
\label{jct:tail-theorem}
Let \(V\in C_c^\infty((1,2))\), \(Y\ge2\), \(R\ge1\), and
\(U>0\) with \(U^2\ge Y\).  Define
\begin{equation}
\begin{split}
 J_V(Y,U,R)={}&\sum_{\mathfrak a_1,\mathfrak a_2}
 m(\mathfrak a_1)m(\mathfrak a_2)
 b(\N\mathfrak a_1/U)b(\N\mathfrak a_2/U)
 V(\N\mathfrak a_1\N\mathfrak a_2/Y)\\
 &\hspace{8mm}\cdot
 \left[1-\chi\left(
 \frac{\cont(\mathfrak a_1)\cont(\mathfrak a_2)}R\right)\right].
 \label{jct:tail}
\end{split}
\end{equation}
Under \eqref{jct:zero-free-input}, for every \(\varepsilon>0\),
\begin{equation}
 J_V(Y,U,R)\ll_{\beta,u,\varepsilon,V,b,\chi}
 Y^{\beta+\varepsilon}R^{-a_\beta}\log(2+U^2/Y).
 \label{jct:tail-general}
\end{equation}
In particular \(\beta=1,u=1/2\) gives the unconditional bound
\(Y^{1+\varepsilon}R^{-3/2}\log(2+U^2/Y)\), while
\(\beta=1,u=1/4\) gives exponent \(-7/4\) on \(R\) under zeta RH.
The constants are uniform over families with common compact support
in \((1,2)\) and bounded finitely many logarithmic-derivative
seminorms, with their number depending on \(\varepsilon\).
The sum is empty for \(R\ge\sqrt{2Y}\).
No coprimality between the two ideals, or between their contents,
is imposed.
\end{theorem}

\begin{theorem}[Specialization of the published zero-free input]
\label{jct:liu-specialization}
Using Liu's zero-free theorem for finite-order Hecke characters,
one may take both \(\beta=B_*\) and \(\rho=B_*\), where
\[
 B_*=0.873345856106718576801184796065\ldots.
\]
The exponents in \eqref{jct:exponents} are then
\[
\begin{array}{c|cc}
 \rho&\gamma_{B_*,\rho/2}&1-2\gamma_{B_*,\rho/2}\\\hline
 1&0.299469269306252\ldots&0.401061461387495\ldots\\
 B_*&0.284992750196395\ldots&0.430014499607209\ldots\\
 1/2&0.249447399928862\ldots&0.501105200142275\ldots.
\end{array}
\]
The middle row uses the published theorem for zeta as well as for
the cubic Hecke function. Its exact cutoff is
\((2B_*-1)/(3B_*)\). The last row additionally assumes classical
zeta RH. Its content
exponent is strictly below \(1/4\), since \(B_*<7/8\).
The first two content exponents are strictly below the elementary
value \(1/3\). The middle exponent is even below the earlier
zeta-RH value \(2/7\), since again \(B_*<7/8\).
The last exponent is below both of those values.
\end{theorem}

Here the cited input is exactly \cite[Theorem~1.2, p.~3;
Theorem~8.29, pp.~78--80]{JCTLiu}, with its fixed-field constants
and its open half-plane. Its long analytic proof is imported,
not reconstructed in this paper. The finite deletion factor \(H\)
has zeros only on \(\Re s=0\), so that theorem gives the first
part of \eqref{jct:zero-free-input}. Its trivial character over
\(\mathbf Q\) gives the zeta assertion, allowing its pole at one.
The tail theorem itself is stated for arbitrary explicitly assumed
\(\beta\) and \(\rho\), independently of this
particular external input. Neither theorem asserts cubic RH.

\paragraph{Weights and the balanced expressions.}
Fix \(r>1\), \(\kappa>4r\), and
\(W\in C_c^\infty((1,2))\) satisfying
\begin{equation}
 \|\Dc^nW\|_\infty
 \le C_WH_W^n n!\{\log(e+n)\}^{rn}\qquad(n\ge0)
 \label{jct:weight}
\end{equation}
for some fixed positive \(C_W,H_W\).
Constants in the balanced assertions may depend on the fixed
parameters \(r\) and \(\kappa\).
For \(X\ge2\), put
\begin{equation}
 U=\sqrt{2X},\qquad
 L_X=\{\log(3X)\}^4\{\log(e+\log(3X))\}^{\kappa}.
 \label{jct:scales}
\end{equation}
Let \(d_{2,K}(\mathfrak c)\) count ordered factorizations of
\(\mathfrak c\) into two integral ideals.  Write
\begin{equation}
\begin{split}
 B_W(\mathfrak c,\mathfrak d_1,\mathfrak d_2;X)={}&
 b(\N\mathfrak d_1/U)b(\N\mathfrak d_2/U)
 W\left(\frac{\N\mathfrak c\N\mathfrak d_1\N\mathfrak d_2}X\right)\\
 &\cdot\left[1-\chi_d\left(
 \frac{L_X\N\mathfrak d_1\N\mathfrak d_2}X\right)\right],
 \label{jct:balanced-weight}
\end{split}
\end{equation}
and define the finite sums
\begin{align}
 R_{\mathrm{soft},W}(X)
 ={}&-\frac12\sum_{\mathfrak c,\mathfrak d_1,\mathfrak d_2}
 \psi(\mathfrak c)d_{2,K}(\mathfrak c)\log\N\mathfrak c\,
 m(\mathfrak d_1)m(\mathfrak d_2)
 B_W(\mathfrak c,\mathfrak d_1,\mathfrak d_2;X),
 \label{jct:rsoft}\\
 R_{\gamma,W}(X)
 ={}-\frac12\sum_{\mathfrak c,\mathfrak d_1,\mathfrak d_2}
 &\psi(\mathfrak c)d_{2,K}(\mathfrak c)\log\N\mathfrak c\,
 m(\mathfrak d_1)m(\mathfrak d_2)
 B_W(\mathfrak c,\mathfrak d_1,\mathfrak d_2;X)\notag\\[-2mm]
 &\hspace{4mm}\cdot\chi\left(
 \frac{\cont(\mathfrak d_1)\cont(\mathfrak d_2)}
 {(X/\N\mathfrak c)^\gamma}\right).
 \label{jct:rjoint}
\end{align}
The shell and cutoffs make explicit upper bounds on these sums
unnecessary; the bounds are part of the next assertion.

\begin{theorem}[A smaller balanced prime convolution]
\label{jct:balanced-theorem}
Let
\[
 T_W(X)=\sum_{\mathfrak a}\Lambda_K(\mathfrak a)\psi(\mathfrak a)
                                  W(\N\mathfrak a/X).
\]
With the preceding definitions, for every \(M>0\),
\begin{equation}
 T_W(X)=R_{\mathrm{soft},W}(X)+O_{W,b,\chi_d,M}(X^{-M}).
 \label{jct:soft-main}
\end{equation}
Take \(\gamma=\gamma_{\beta,u}\) from \eqref{jct:exponents}.
Under the stated zero-free input, for every \(\varepsilon>0\),
\begin{equation}
 T_W(X)=R_{\gamma,W}(X)+O_{\beta,\rho,W,b,\chi,\chi_d,\varepsilon}
                                      (X^{1/2+\varepsilon}).
 \label{jct:joint-main}
\end{equation}
Every nonzero summand in either balanced sum satisfies
\begin{equation}
 \N\mathfrak c<2L_X,\qquad
 \frac{U}{4L_X}<\N\mathfrak d_i<2U\quad(i=1,2).
 \label{jct:balanced-support}
\end{equation}
In \(R_{\gamma,W}\), write
\(\mathfrak d_i=(r_i)\mathfrak b_i\) with primitive
\(\mathfrak b_i\), and put \(q_i=\N\mathfrak b_i\).
Then
\begin{equation}
 q_1q_2>\frac14(X/\N\mathfrak c)^{1-2\gamma}.
 \label{jct:primitive-threshold}
\end{equation}
The choices \((\beta,\rho)=(1,1)\) and \((1,1/2)\) give
exponents \(1/3\) and \(3/7\) in this primitive-product bound.
For \((\beta,\rho)=(1,B_*)\), the cutoff is \(1/(4-B_*)\)
and the primitive exponent is \((2-B_*)/(4-B_*)\).
The exponents with \(\beta=B_*\) are
listed in Theorem~\ref{jct:liu-specialization}. For every nonzero
nonnegative admissible \(W\), infinitely many scales have nonzero
summands removed by the content cutoff. More precisely, for each
of the three rows in Theorem~\ref{jct:liu-specialization}, infinitely
many scales have a nonzero
summand that is unchanged by the \(\beta=1\) cutoff and suppressed
by the \(\beta=B_*\) cutoff. The \(\rho=B_*\) cutoff also suppresses
nonzero summands left unchanged by the \(\beta=B_*,\rho=1\)
cutoff. Thus each refinement deletes a strictly larger region.
No estimate for \(R_{\gamma,W}\) itself
is asserted.
\end{theorem}

\begin{theorem}[The exact remaining square-root criterion]
\label{jct:rh-theorem}
Fix \(r>1\), \(\kappa>4r\), and the three cutoffs above.
Assume \eqref{jct:zero-free-input}, retain the specified choice of
\(u\), and take \(\gamma=\gamma_{\beta,u}\). Define
\[
 t(p)=\#\{z\bmod p:z^3=269\}-1,\qquad
 \Theta_W(X)=\sum_{\substack{p\ne3,269\\p\equiv1\ (3)}}
 t(p)\log p\,W(p/X).
\]
The following are equivalent:
\begin{enumerate}
\item RH for the primitive cubic Hecke function \(L(s)\).
\item For every nonnegative \(W\) satisfying \eqref{jct:weight}
and every \(\varepsilon>0\),
\(R_{\gamma,W}(X)=O_{W,\varepsilon}(X^{1/2+\varepsilon})\).
\item For every such \(W\) and every \(\varepsilon>0\),
\(\Theta_W(X)=O_{W,\varepsilon}(X^{1/2+\varepsilon})\).
\end{enumerate}
Nonzero nonnegative admissible weights exist in every open
subinterval of \((1,2)\).  Central and multiple zeros are allowed;
the criterion is quantified over all the indicated weights,
not over one selected weight.
\end{theorem}

\section{Proof overview}
The squarefree ideal coefficients have a two-variable Euler product,
recording their norm and rational content together. Its local factor
is \(1-t_pp^{-s}+\eta(p)p^{-v}\). Extracting the inverse cubic
Hecke function in \(s\) and an elementary rational series in \(v\)
first leaves a mixed term at \(s+v\). Extracting this additional
Hecke factor leaves a normally
convergent correction throughout \(\Re s>\beta\), \(\Re v>u\).
The zero-free input therefore permits moving the norm lines to just
right of \(\beta\) and the content lines just right of \(u\),
even when \(\beta+u<1\). This gives the exponent \(2\beta-u\).

A Mellin transform for the two norm caps and a third Mellin variable
for the product of contents turn that factorization into the tail
estimate. The contours are moved one at a time within the proved
domain. The cap powers cancel on the final lines, leaving only the
explicit logarithm of their aspect ratio. A separate rational-content
summation also proves the \(\beta=1\) cases without invoking the
bivariate factorization.
The smooth cap is then built into the Heath--Brown identity itself.
The inverse tail starts above \(U\), so its square cannot meet
the shell of norm at most \(2X\).  Entireness and the degree-two
Hecke functional equation remove the first term of the identity.
For its second term, a contour estimate uniform in the shift
and the fixed derivative condition restrict the direct ideal to
norm \(O(L_X)\).  Applying the joint-content estimate for each
such direct ideal gives \eqref{jct:joint-main}.
Finally, ordinary prime-power counting and Mellin continuation
give the all-weight RH criterion.

\section{Elementary rational-content factorization}
Put \(\eta=\chi_{-3}\).  The Eisenstein integer ring is norm
Euclidean, hence a PID, and its norm is
\(a^2-ab+b^2\ge(a^2+b^2)/2\).
There are six units.  Consequently the number of ideals of norm
at most \(Z\) is \(O(Z)\), by counting their generators in the
corresponding ellipse.  The same estimate holds after deleting
any fixed finite set of primes.

At a split rational prime \(p\), write
\((p)=\mathfrak p\overline{\mathfrak p}\).
The content valuation of an ideal is the minimum of its two
prime-ideal valuations.  At an inert prime, \((p)\) itself is
prime and its valuation is the content valuation.
The ramified prime 3 is deleted.  These local descriptions prove
the unique decomposition
\(\mathfrak a=(r)\mathfrak b\) with primitive \(\mathfrak b\).

A primitive ideal with nonzero \(\mu_K\) has no inert prime
factor and has at most one of the two primes above each split
prime, with exponent one.  Its norm is therefore squarefree and
supported on split primes.  For such an ideal, the ideal
\((r)\mathfrak b\) is squarefree exactly when \(r\) is squarefree
and \((r,\N\mathfrak b)=1\).
At a split prime \(\mu_K((p))=1=\eta(p)\); at an inert prime
\(\mu_K((p))=-1=\eta(p)\).  Using \(\psi((r))=1\) gives the
exact identity, including its zero cases,
\begin{equation}
 m((r)\mathfrak b)=\mu(r)^2\eta(r)m(\mathfrak b)
                         \mathbf1_{(r,\N\mathfrak b)=1}.
 \label{jct:content-factorization}
\end{equation}
Define
\begin{equation}
 h(q)=\sum_{\substack{\N\mathfrak b=q\\\cont(\mathfrak b)=1}}
 m(\mathfrak b),\qquad
 \sum_{q\le Z}|h(q)|\ll Z.
 \label{jct:primitive-count}
\end{equation}
The inequality follows by taking absolute values inside the sum
and applying the ideal count.  It does not require cancellation.
Partial summation twice yields
\begin{align}
 \sum_{q_1q_2\le Q}|h(q_1)h(q_2)|&\ll Q\log(2Q),
 \label{jct:pair-count}\\
 \sum_{q_1q_2\le Q}\frac{|h(q_1)h(q_2)|}{(q_1q_2)^v}
 &\ll_v Q^{1-v}\log(2Q)\quad(0\le v<1).
 \label{jct:weighted-pair-count}
\end{align}
For the first bound, sum \(q_2\) first and use
\(\sum_{q\le Q}|h(q)|/q\ll\log(2Q)\); the second follows
by partial summation in the product.
Also, counting ordered ideal factorizations gives
\begin{equation}
 \sum_{\N\mathfrak c\le Z}d_{2,K}(\mathfrak c)
 \ll Z\log(2Z).
 \label{jct:divisor-count}
\end{equation}

When \eqref{jct:content-factorization} is applied to two ideals,
the only coprimality conditions are \((r_i,q_i)=1\) individually.
The two contents and the two primitive norms may share primes;
there is no condition on \((r_1,q_2)\) or \((r_2,q_1)\).
In particular \(\mu(r_1)^2\mu(r_2)^2\) is not
\(\mu(r_1r_2)^2\).
Nor is \(r_1r_2\) necessarily the content of the product ideal:
conjugate prime factors in the primitive ideals may increase
that content.

\section{The rational series and all finite corrections}
For a primitive norm \(q\) in the support of \(h\), set
\[
 a_q(r)=\mu(r)^2\eta(r)\mathbf1_{(r,3\cdot269q)=1},\qquad
 A_q(s)=\sum_{r\ge1}a_q(r)r^{-s}.
\]
The local identity \(1+z=(1-z^2)/(1-z)\) gives, initially for
\(\Re s>1\),
\begin{equation}
 A_q(s)=\frac{L(s,\eta)}
 {\zeta(2s)(1-3^{-2s})(1-269^{-s})}
 \prod_{p\mid q}(1+p^{-s})^{-1}.
 \label{jct:rational-series}
\end{equation}
At 3 the factor from \(1/\zeta(2s)\) must be removed, since
\(\eta(3)=0\).  At 269 one has \(\eta(269)=-1\), and
deleting its original factor requires \((1-269^{-s})^{-1}\).
Every prime dividing \(q\) is split, giving the final product.
Thus the deletion at 269 is a separate correction even when
\(q=1\).

For every \(\sigma_0,\nu>0\),
\begin{equation}
 \prod_{p\mid q}|1+p^{-s}|^{-1}
 \le\prod_{p\mid q}(1-p^{-\sigma_0})^{-1}
 \ll_{\sigma_0,\nu}q^\nu\qquad(\Re s\ge\sigma_0).
 \label{jct:q-uniformity}
\end{equation}
Indeed, outside a fixed finite set of primes each factor is at
most \(p^\nu\).  This proves uniformity even when the number of
prime divisors of \(q\) grows.
Bounded partial sums of \(\eta\) give holomorphy and polynomial
vertical growth of \(L(s,\eta)\) on every closed strip in
\(\Re s>0\), by partial summation.
Unconditionally, \(1/\zeta(2s)\) is absolutely convergent for
\(\Re s>1/2\).

\begin{lemma}[Reciprocal growth in a zero-free half-plane]
\label{jct:reciprocal-growth}
Let \(f\) be holomorphic and nonvanishing in \(\Re z>a\), with
polynomial growth on every bounded closed vertical strip there.
Suppose that on one fixed line \(\Re z=b>a\),
\(|f(b+it)|\ge c>0\) for all real \(t\). Then \(1/f\) has
polynomial growth on every bounded closed strip strictly inside
\(\Re z>a\).
\end{lemma}
\begin{proof}
Fix a target real-part interval \([a+\delta,B]\) and put
\(T=3+|t|\). Choose a fixed-width rectangle around the horizontal
segments joining \(b+it\) to this interval at height \(t\),
with fixed positive margins and left edge to the right of \(a\).
Polynomial strip growth gives
\(\log|f(z)|\le C_0+A\log T\) throughout that rectangle.
Increase \(C_0\) so that
\[
 h_t(z)=C_0+A\log T-\log|f(z)|
\]
is positive and harmonic there. Harmonicity follows from
nonvanishing; no global choice of a complex logarithm is needed.
Cover each horizontal segment by a uniformly bounded number of
disks of fixed positive radius, with successive centers at most
half a radius apart. The Poisson formula for a positive harmonic
function bounds its value within half a disk radius by three times
its value at the center. Successive applications therefore give
\[
 h_t(\sigma+it)\le C_Hh_t(b+it)\ll\log T,
 \qquad a+\delta\le\sigma\le B.
\]
The last inequality uses the assumed lower bound at \(b+it\).
Thus \(\log|f(\sigma+it)|\ge-C\log T\), proving
\(|1/f(\sigma+it)|\ll T^C\), with constants independent of
\(t\) and \(\sigma\) in the target strip.
\end{proof}

Apply the lemma to \(f(z)=(z-1)\zeta(z)\) in \(\Re z>\rho\).
The pole at one is removable in \(f\), with nonzero value there;
standard zeta strip growth makes \(f\) polynomially bounded.
On \(\Re z=2\), the absolutely convergent Euler product bounds
\(1/\zeta(z)\), and \(|z-1|\ge1\), giving the required lower
bound for \(f\). Hence
\(1/\zeta(z)=(z-1)/f(z)\) has polynomial strip growth; it is
holomorphic at one, where it vanishes. The zeta zero-free input
therefore makes \(1/\zeta(2s)\) holomorphic with polynomial
vertical growth on every closed strip to the right of \(u=\rho/2\).

Recall that \(u=\rho/2\).
For every fixed small \(\delta>0\) and \(\nu>0\),
\begin{equation}
 |A_q(s)|\ll_{\delta,\nu}q^\nu(1+|\Im s|)^{B_\delta},
 \qquad u+\delta\le\Re s\le2.
 \label{jct:series-strip}
\end{equation}
There are no poles in this strip.  The elementary denominators
in \eqref{jct:rational-series} have their zeros on \(\Re s=0\).

\section{Keeping norm and content in the same Euler product}
\label{jct:zero-free-proof}
This section proves the general tail theorem.  Its extra input is
exactly the two zero-free half-planes stated in Section~1. The
specialization \(\beta=\rho=1\) requires no additional theorem.

\paragraph{The two complex variables.}
For \(\Re s>1\) and \(\Re v>1\), define the absolutely convergent series
\begin{equation}
 \mathcal F(s,v)=\sum_{\mathfrak a}
 m(\mathfrak a)(\N\mathfrak a)^{-s}
 \cont(\mathfrak a)^{2s-v}.
 \label{jct:bivariate}
\end{equation}
The powers of positive integers use their real logarithms.  Let
\(\eta=\chi_{-3}\), and at a rational prime \(p\ne3,269\) put
\[
 t_p=\begin{cases}\psi(\mathfrak p)+\psi(\overline{\mathfrak p}),
                   &p\text{ split in }K,\\0,&p\text{ inert},\end{cases}
 \qquad D_p(s)=1-t_pp^{-s}+\eta(p)p^{-2s}.
\]
Thus \(P(s)^{-1}=\prod_{p\ne3,269}D_p(s)\).
At a split prime the four squarefree ideal possibilities are
\(1,\mathfrak p,\overline{\mathfrak p},(p)\).
Their contributions to \eqref{jct:bivariate} are
\(1,-\psi(\mathfrak p)p^{-s},
-\psi(\overline{\mathfrak p})p^{-s},p^{-v}\).
At an inert prime the two possibilities contribute \(1,-p^{-v}\).
Consequently the local factor is, in both cases,
\begin{equation}
 1-t_pp^{-s}+\eta(p)p^{-v}.
 \label{jct:bivariate-local}
\end{equation}
This computation includes ideals of content divisible by primes
in the other ideal; no cross-coprimality is introduced.

Put
\begin{equation}
 A(v)=A_1(v)=\frac{L(v,\eta)}
 {\zeta(2v)(1-3^{-2v})(1-269^{-v})},\qquad
 E(s,v)=\prod_{p\ne3,269}
 \frac{1-t_pp^{-s}+\eta(p)p^{-v}}
 {(1+\eta(p)p^{-v})D_p(s)}.
 \label{jct:bivariate-factor}
\end{equation}
The factors at both deleted rational primes are displayed explicitly.
The local computation proves
\begin{equation}
 \mathcal F(s,v)=A(v)P(s)^{-1}E(s,v)
 \label{jct:bivariate-continuation}
\end{equation}
first in the absolutely convergent region.

The mixed term in \(E\) can be extracted as another Hecke factor.
Write \(x=p^{-s}\), \(y=p^{-v}\), \(e=\eta(p)\), and \(t=t_p\).
At a split prime \(e=1\), and at an inert prime \(t=0\), so
\(et=t\) in both cases. Define
\begin{equation}
 E_{2,p}(s,v)=
 \frac{(1-tx+ey)(1-txy+ex^2y^2)}
 {(1+ey)(1-tx+ex^2)}.
 \label{jct:secondary-local}
\end{equation}
Local multiplication gives, initially in the absolute region,
\begin{equation}
 E(s,v)=P(s+v)E_2(s,v),\qquad
 E_2(s,v)=\prod_{p\ne3,269}E_{2,p}(s,v).
 \label{jct:secondary-factor}
\end{equation}
All local factors at the deleted primes remain omitted in every
product. Subtracting the denominator polynomial in
\eqref{jct:secondary-local} from its numerator gives exactly
\[
 \begin{split}
 -txy^2+x^2\{(t^2-1)y+ey^2+y^3-e\}-tx^3y^2.
 \end{split}
\]
Since \(|t|\le2\) and \(|x|,|y|<1\), this is
\(O(|x|^2+|x||y|^2)\).
The inverse roots of \(D_p\) have modulus one. Thus on
\(\Re s\ge\sigma_0>0\), \(\Re v\ge\nu_0>0\), the denominator
has absolute value at least
\[
 (1-p^{-\nu_0})(1-p^{-\sigma_0})^2,
\]
with a positive lower bound independent of the imaginary parts.
Consequently
\[
 E_{2,p}(s,v)=1+O\bigl(p^{-2\Re s}+p^{-\Re(s+2v)}\bigr).
\]
The product \(E_2\) converges normally, and is uniformly bounded
on closed inner regions, when
\begin{equation}
 \Re s>\tfrac12,\qquad\Re v>0,\qquad\Re(s+2v)>1.
 \label{jct:secondary-domain}
\end{equation}
No nonvanishing assertion for \(E_2\) is needed. Because
\(\beta>1/2\) and \(u\ge1/4\), every point with
\(\Re s>\beta\), \(\Re v>u\) lies strictly inside this domain:
in particular \(\Re(s+2v)>\beta+2u>1\).

\paragraph{Vertical bounds without a simple-zero hypothesis.}
Apply Lemma~\ref{jct:reciprocal-growth} to \(P\) in
\(\Re s>\beta\). Its strip growth is part of the stated Hecke
background, and its absolutely convergent reciprocal Euler product
on \(\Re s=2\) gives the required lower bound there.
The zero-free input consequently yields
\[
 |P(s)^{-1}|\ll_\delta(1+|\Im s|)^{C_\delta}
\]
on every fixed bounded strip with \(\Re s\ge\beta+\delta\).
This argument requires neither a zero-count estimate nor any
restriction on zero multiplicities.

The rational-series argument above gives the corresponding
polynomial bound for \(A(v)\) on every closed strip in
\(\Re v>u=\rho/2\), under the stated zeta zero-free input.
The pole of \(\zeta\) becomes a zero of its reciprocal and causes
no singularity. The additional factor \(P(s+v)\) in
\eqref{jct:secondary-factor} is entire and polynomially bounded
on the fixed strips in question. Its zeros are harmless because
it is not inverted.
Equations~\eqref{jct:bivariate-continuation} and
\eqref{jct:secondary-factor} therefore continue \(\mathcal F\)
holomorphically, with polynomial vertical growth, to
\begin{equation}
 \Re s>\beta,\qquad \Re v>u.
 \label{jct:tube}
\end{equation}
The finite corrections in \eqref{jct:bivariate-factor} have no
zeros in their denominators there.

\paragraph{The two norm caps and the content cutoff.}
Put \(q=U^2/Y\ge1\) and
\[
 \Psi_q(t_1,t_2)=b(t_1)b(t_2)V(qt_1t_2),\qquad
 K(t)=1-\chi(t).
\]
The support of \(\Psi_q\) lies in
\[
 0<t_i\le2,\qquad q^{-1}<t_1t_2<2q^{-1}.
\]
It is compact in the positive quadrant for each fixed \(q\),
since \(t_i>1/(2q)\).  Define
\[
 \widehat\Psi_q(s_1,s_2)=\iint
 \Psi_q(t_1,t_2)t_1^{s_1}t_2^{s_2}
 \frac{dt_1}{t_1}\frac{dt_2}{t_2}.
\]
For fixed \(\sigma>0\), logarithmic coordinates show that the
support has area \(O(\log(2+q))\) and that
\((t_1t_2)^\sigma\ll q^{-\sigma}\).
Each fixed logarithmic derivative of \(\Psi_q\) is bounded
uniformly in \(q\): differentiating \(V(qt_1t_2)\) acts on
its argument.  Repeated integration by parts against the full
Mellin parameters consequently proves, for every integer \(J\ge0\),
\begin{equation}
 |\widehat\Psi_q(\sigma+it_1,\sigma+it_2)|
 \ll_{\sigma,J}q^{-\sigma}\log(2+q)
 (1+|t_1|)^{-J}(1+|t_2|)^{-J}.
 \label{jct:norm-caps}
\end{equation}
For real parts \(\sigma_i\in[\sigma,2]\), the additional
factors \(t_i^{\sigma_i-\sigma}\) are bounded by fixed constants,
so the same estimate controls the intervening contours.
All constants use only finitely many of the stated weight seminorms.

For \(\Re z>0\), define the Mellin transform with the opposite sign
\begin{equation}
 \widetilde K(z)=\int_0^\infty K(t)t^{-z}\frac{dt}{t}
 =\frac1z\int_0^\infty (tK'(t))t^{-z}\frac{dt}{t}.
 \label{jct:cutoff-mellin}
\end{equation}
The boundary terms vanish because \(K=0\) through one and
\(K=1\) from two onward.  Its derivative is smooth and supported
in \([1,2]\).  Hence \(\widetilde K\) is holomorphic for
\(\Re z>0\), and decreases faster than every inverse power of
\(|\Im z|\) on each fixed positive real line.  Mellin inversion gives
\(K(t)=(2\pi i)^{-1}\int_{(\alpha)}\widetilde K(z)t^z\,dz\)
for every fixed \(\alpha>0\).

\paragraph{Successive contour shifts.}
Choose a sufficiently small \(\delta>0\), and put
\[
 \sigma=\beta+\delta,\qquad \nu=u+2\delta,\qquad
 \alpha=2\sigma-\nu=a_\beta>0.
\]
Take \(\delta\) small enough that \(\sigma<2\).
Use the \(z\)-line of real part \(\alpha\) from the outset.
On the initial two norm lines \(\Re s_i=2\), the two values
\(v_i=2s_i-z\) have real part \(4-\alpha>1\).
Absolute convergence, three Mellin inversions and Fubini give
\begin{equation}
\begin{split}
 J_V(Y,U,R)=\frac1{(2\pi i)^3}
 \int_{(\alpha)}\int_{(2)}\int_{(2)}&
 \widetilde K(z)R^{-z}\widehat\Psi_q(s_1,s_2)
 U^{s_1+s_2}\\
 &\cdot\mathcal F(s_1,2s_1-z)
 \mathcal F(s_2,2s_2-z)\,ds_2\,ds_1\,dz.
\end{split}
 \label{jct:triple-inversion}
\end{equation}
Move first the \(s_1\)-line, and then the \(s_2\)-line, to real
part \(\sigma\), keeping the \(z\)-line fixed.  Throughout either
shift, with \(\sigma\le\Re s_i\le2\),
\[
 \Re s_i\ge\sigma>\beta,\qquad
 \Re(2s_i-z)\ge\nu>u.
\]
Thus both factors stay inside \eqref{jct:tube}; no pole is crossed.
Their polynomial bounds depend only polynomially on
\(1+|\Im s_1|+|\Im s_2|+|\Im z|\).
Choose the decay orders in \eqref{jct:norm-caps} and
\eqref{jct:cutoff-mellin} larger than these fixed powers plus four.
The resulting integrable majorant justifies the successive
horizontal limits and all interchanges, including integration in
\(z\).  This is not a simultaneous contour deformation assumption.

On the final lines the scale factors are exactly
\[
 U^{2\sigma}q^{-\sigma}R^{-\alpha}
 =Y^\sigma R^{-\alpha}.
\]
The integral is consequently
\[
 J_V(Y,U,R)\ll Y^{\beta+\delta}
 R^{-a_\beta}\log(2+U^2/Y).
\]
Since \(R\ge1\), take \(\delta\le\varepsilon\) to obtain the
claimed bound.  A nonzero term has content product greater than
\(R\) and norm product at least the square of that content product,
while the shell has norm product less than \(2Y\).
This proves the empty case and completes the general tail theorem,
including its bounded-family uniformity.


\section{Separate rational summation}
This independent argument proves the two elementary specializations
of Theorem~\ref{jct:tail-theorem}. The general zero-free improvement
was proved in Section~\ref{jct:zero-free-proof}.

A nonzero term of \eqref{jct:tail} has \(r_1r_2>R\) and
\(q_1q_2r_1^2r_2^2<2Y\).  Thus
\begin{equation}
 q_1q_2<Q_0:=2Y/R^2.
 \label{jct:q-range}
\end{equation}
This also proves the empty case in Theorem~\ref{jct:tail-theorem}.
In the remaining case, \eqref{jct:content-factorization} gives
\begin{equation}
 J_V(Y,U,R)=\sum_{q_1q_2<Q_0}h(q_1)h(q_2)S_{q_1,q_2},
 \label{jct:decomposition}
\end{equation}
where
\begin{equation}
\begin{split}
 S_{q_1,q_2}=\sum_{r_1,r_2\ge1}&a_{q_1}(r_1)a_{q_2}(r_2)
 b(q_1r_1^2/U)b(q_2r_2^2/U)\\
 &\cdot V(q_1q_2r_1^2r_2^2/Y)[1-\chi(r_1r_2/R)].
 \label{jct:inner-content-sum}
\end{split}
\end{equation}
Put
\[
 z=U^2/Y\ge1,\quad Z_i=(U/q_i)^{1/2},\quad
 \alpha=Z_1Z_2/R,
\]
and define
\begin{equation}
 \Phi_{z,\alpha}(t_1,t_2)=b(t_1^2)b(t_2^2)
 V(z(t_1t_2)^2)[1-\chi(\alpha t_1t_2)].
 \label{jct:two-weight}
\end{equation}
Then the weight in \eqref{jct:inner-content-sum} is
\(\Phi_{z,\alpha}(r_1/Z_1,r_2/Z_2)\).
Every fixed mixed derivative in
\(\Dc_i=t_i\partial/\partial t_i\) is uniformly bounded in
\(z\ge1\) and \(\alpha>0\), in terms of finitely many fixed
seminorms of the three weights.  For example,
\(\Dc_i V(z(t_1t_2)^2)=2(\Dc V)(z(t_1t_2)^2)\).
Logarithmic differentiation of the moving cutoff similarly
acts on its argument, without an uncontrolled scale factor.

The weight and all its derivatives are supported in
\begin{equation}
 0<t_1,t_2\le\sqrt2,\qquad
 z^{-1/2}<t_1t_2<\sqrt2\,z^{-1/2}.
 \label{jct:mellin-region}
\end{equation}
This is contained in a compact subset of the positive quadrant
for each fixed \(z\), since \(t_i\ge1/\sqrt{2z}\).
For fixed \(\sigma>0\),
\begin{equation}
 \iint_{\eqref{jct:mellin-region}}
 (t_1t_2)^\sigma\frac{dt_1}{t_1}\frac{dt_2}{t_2}
 \ll_\sigma z^{-\sigma/2}\log(2+z).
 \label{jct:aspect}
\end{equation}
To verify the aspect factor, set \(v_i=\log t_i\).
The sum \(v_1+v_2\) lies in a fixed-length interval at
\(-\tfrac12\log z\).  For each such sum, \(v_1\) ranges
over an interval of length \(O(\log(2+z))\), and
\(\exp(\sigma(v_1+v_2))\ll z^{-\sigma/2}\).

Integrating by parts in both variables against the full complex
Mellin parameters now gives, for every integer \(J\ge0\),
\begin{equation}
 |\widehat\Phi(\sigma+it_1,\sigma+it_2)|
 \ll_{\sigma,J}z^{-\sigma/2}\log(2+z)
 (1+|t_1|)^{-J}(1+|t_2|)^{-J}.
 \label{jct:double-decay}
\end{equation}
For intermediate contours with real parts \(\sigma_i\in[\sigma,2]\),
the extra factors \(t_i^{\sigma_i-\sigma}\) are bounded because
\(t_i\le\sqrt2\).  The same estimate with a fixed additional
constant therefore controls these contours as well.
All constants involve only finitely many of the stated
seminorms, uniformly in \(Y,U,R,q_1,q_2\).

Initially on two lines of real part two, absolute convergence
and Mellin inversion give
\begin{equation}
 S_{q_1,q_2}=\frac1{(2\pi i)^2}\int_{(2)}\int_{(2)}
 A_{q_1}(s_1)A_{q_2}(s_2)\widehat\Phi(s_1,s_2)
 Z_1^{s_1}Z_2^{s_2}\,ds_2\,ds_1.
 \label{jct:double-inversion}
\end{equation}
Move the contours successively, not as an unproved simultaneous
limit, to \(\Re s_i=\sigma=u+\delta\).
Taking \(J>B_\delta+2\) in \eqref{jct:double-decay} justifies
the horizontal limits and all Fubini interchanges.  For each
fixed set of scale parameters the additional factors on the
intermediate contours are finite; their bounds suffice for
these limits.  There are no residues.  On the final contours
the scale dependence cancels exactly:
\begin{equation}
 (Z_1Z_2)^\sigma z^{-\sigma/2}
 =\left(\frac{Y}{q_1q_2}\right)^{\sigma/2}.
 \label{jct:scale-cancellation}
\end{equation}
Hence
\begin{equation}
 |S_{q_1,q_2}|\ll_{\delta,\nu,V,b,\chi}
 (q_1q_2)^\nu
 \left(\frac{Y}{q_1q_2}\right)^{(u+\delta)/2}
 \log(2+U^2/Y).
 \label{jct:inner-bound}
\end{equation}

Apply \eqref{jct:weighted-pair-count} with \(v=u/2\), using
\(q_1q_2<2Y\) to absorb its positive \(\nu\)-power and
discarding the extra negative \(\delta/2\)-power.  We obtain
\begin{align}
 |J_V|&\ll Y^{\nu+\delta/2+u/2}
 Q_0^{1-u/2}\log(2Q_0)\log(2+U^2/Y)\notag\\
 &\ll Y^{1+\nu+\delta/2}R^{u-2}
 \log(2Y)\log(2+U^2/Y).
 \label{jct:tail-before-epsilon}
\end{align}
Choose \(\nu,\delta\) sufficiently small for the prescribed
\(\varepsilon\) and absorb \(\log(2Y)\).
This proves the two \(\beta=1\) tail bounds and their bounded-family
uniformity.
The two rational variables do not double the exponent of \(R\):
the estimated scale is their product.  Its larger allowable
range is what makes the application useful.

\section{Direct contour bounds with uniform constants}
For \(j=1,2\), let
\begin{equation}
 F_j(s)=-\frac{(P^j)'(s)}j=-P'(s)P(s)^{j-1}
 =\sum_{\mathfrak c}\psi(\mathfrak c)d_{j,K}(\mathfrak c)
 \frac{\log\N\mathfrak c}{j}(\N\mathfrak c)^{-s}
 \label{jct:direct-series}
\end{equation}
in the absolutely convergent half-plane, where \(d_{1,K}=1\).
Write \(A_{j,W}(Y)\) for its coefficients summed against
\(W(\N\mathfrak c/Y)\).
Set \(\mathcal B=\max(1,3N/(4\pi^2))\).
We prove the bound, uniform in every integer \(A\ge1\),
\begin{equation}
 |A_{j,W}(Y)|\le C_j\mathcal B^{jA}Y^{-A}
 \|\Dc^{2jA+j+3}W\|_\infty\qquad(Y\ge1).
 \label{jct:uniform-direct}
\end{equation}

Put \(G(s)=\Gamma(1-s)/\Gamma(s)\) and
\(R_A(t)=(A^2+t^2)^{1/2}\).  Gamma recurrence and conjugation
give the exact formula
\begin{equation}
 |G(-A+it)|=|t|\prod_{v=1}^A(v^2+t^2)
 \le R_A(t)^{2A+1}.
 \label{jct:gamma}
\end{equation}
Reflection gives \(G(s)=\pi^{-1}\sin(\pi s)\Gamma(1-s)^2\).
For \(t\ne0\), differentiate this product and use
\[
 |t\coth(\pi t)|\ll1+|t|,\qquad
 |\operatorname{digamma}(1+A-it)|\ll\log(3+R_A(t)).
\]
The digamma bound is uniform in the displayed right half-plane;
the required gamma formulas and asymptotic are
\cite[5.5.1, 5.5.3, 5.11.2]{JCTDLMF}.
After multiplication, the expression has a removable limit at
\(t=0\), and the result is
\begin{equation}
 |G'(-A+it)|\ll R_A(t)^{2A+1}\log(3+R_A(t)).
 \label{jct:gamma-prime}
\end{equation}
This calculation does not divide by the trivial zero at \(t=0\).
On \(\Re z\ge2\), the Euler series bound both
\(L^\vee(z)\) and its derivative uniformly.
Also \(|H(-A+it)|+|H'(-A+it)|\ll3^A\).
Differentiating \eqref{jct:fe} therefore gives
\[
 |P(-A+it)|\le C\mathcal B^A R_A(t)^{2A+1},\qquad
 |P'(-A+it)|\le C\mathcal B^A R_A(t)^{2A+1}
                                      \log(3+R_A(t)).
\]
The constant includes the fixed conductor logarithm, and contains
no unrecorded factorial in \(A\).  Consequently
\begin{equation}
 |F_j(-A+it)|\le C_j\mathcal B^{jA}
 R_A(t)^{2jA+j}\log(3+R_A(t)).
 \label{jct:direct-strip}
\end{equation}

With \(w(v)=W(e^v)\), integration by parts against the full
complex parameter gives
\begin{equation}
 |\widehat W(-A+it)|\le(\log2)
 \|\Dc^\ell W\|_\infty|-A+it|^{-\ell}.
 \label{jct:weight-mellin}
\end{equation}
Indeed \(w\) is supported in \((0,\log2)\), so the factor
\(e^{-Av}\) is at most one; no factor \(A^\ell\) is introduced.
For each fixed \(A\), entireness, fixed-strip polynomial bounds,
and rapid Mellin decay justify moving the inversion contour for
\(F_j\widehat W Y^s\) from real part two to \(-A\).
The horizontal limits are taken with \(A\) fixed.
Use \(\ell=2jA+j+3\) on the final line.
The remaining integral is bounded uniformly in \(A\ge1\), since
\[
 \int_{\mathbf R}(A^2+t^2)^{-3/2}
 \log(3+\sqrt{A^2+t^2})\,dt\ll1.
\]
This proves \eqref{jct:uniform-direct}; only this already uniform
final estimate is subsequently optimized in \(A\).

Under \eqref{jct:weight}, the inequality \(n!\le n^n\) converts
\eqref{jct:uniform-direct} into an upper bound of the form
\[
 C_1(A+1)^{C_2}\{\log(e+A)\}^{C_3}
 \left\{\frac{C_4(A+1)^{2j}\{\log(e+A)\}^{2jr}}Y\right\}^A.
\]
Take
\(A=\lfloor\delta Y^{1/(2j)}/\{\log(3Y)\}^r\rfloor\)
with a sufficiently small fixed positive \(\delta\).
The bracket is bounded away from one for large \(Y\), and its
exponential absorbs the polynomial factors.  Enlarging the
constant on a bounded initial range gives
\begin{equation}
 |A_{j,W}(Y)|\ll_{j,W}
 \exp\left\{-c_{j,W}\frac{Y^{1/(2j)}}{\{\log(3Y)\}^r}\right\}.
 \label{jct:direct-decay}
\end{equation}
The constants are uniform over families with fixed bounds on
\(C_W,H_W,r\), and not over arbitrarily varying weights.

\paragraph{Existence of localized nonnegative weights.}
Given an open interval \(J\subset(1,2)\), choose
\(a_n=c/(n\{\log(e+n)\}^r)\), with \(c>0\) small enough
that \(\sum a_n<|\log J|\).  Convolve the probability densities
of the intervals \([-a_n/2,a_n/2]\).
Their infinite convolution exists, since the sum of the widths
is finite.  Differentiating the first \(m\) interval measures
and bounding the next density gives
\[
 \|w^{(m)}\|_\infty\le\frac{2^m}{a_1\cdots a_{m+1}}.
\]
Keeping one further density shows each distributional derivative
is continuous: the convolution of two interval densities is
continuous, and convolution with the remaining probability
measure preserves continuity.  Thus \(w\) is smooth,
nonnegative, and has integral one.  The product bound is at most
\(CH^m m!\{\log(e+m)\}^{rm}\), with the extra factors absorbed
in \(H^m\).  Translate its support strictly inside \(\log J\)
and set \(W(t)=w(\log t)\).  This proves the claimed localization
and \eqref{jct:weight}.

\section{Exact smooth capping and the joint deletion}
Define the finite inverse series
\[
 M_{U,b}(s)=\sum_{\mathfrak d}m(\mathfrak d)
 b(\N\mathfrak d/U)(\N\mathfrak d)^{-s}.
\]
Because the cap is identically one through \(U\), the coefficients
of \(D-M_{U,b}\) vanish at every norm at most \(U\).
The classical algebra, for this particular finite series, is
\begin{equation}
 D=2M_{U,b}-PM_{U,b}^2+P(D-M_{U,b})^2.
 \label{jct:identity}
\end{equation}
It follows by expanding the right side and using \(PD=1\),
either in the absolutely convergent half-plane or as an ideal
convolution identity.  This is not a smoothing of an old
coefficient equality: \eqref{jct:identity} is exact for the smooth
cap itself.
The last term has support strictly above \(U^2=2X\), even after
multiplication by \(-P'\), since multiplying ideals cannot
decrease their norms.  Extracting the shell gives exactly
\begin{equation}
\begin{split}
 T_W(X)={}&2\sum_{\mathfrak d}m(\mathfrak d)b(\N\mathfrak d/U)
 A_{1,W}(X/\N\mathfrak d)\\
 &-\sum_{\mathfrak d_1,\mathfrak d_2}
 m(\mathfrak d_1)m(\mathfrak d_2)
 b(\N\mathfrak d_1/U)b(\N\mathfrak d_2/U)
 A_{2,W}(X/(\N\mathfrak d_1\N\mathfrak d_2)).
 \label{jct:exact-shell}
\end{split}
\end{equation}
The coefficient \(1/2\) in the second direct sum comes from
\(-PP'=-(P^2)'/2\); the minus sign is the second coefficient
of \eqref{jct:identity}.  No derivative of \(M_{U,b}\) is taken.

The first sum has \(O(U)\) ideals and direct scale at least
\(X/(2U)\gg\sqrt X\).  A sufficiently large fixed \(A\) in
\eqref{jct:uniform-direct} makes it \(O_M(X^{-M})\) for every
\(M\), even for arbitrary fixed smooth \(W\).
Split the second sum using
\(\chi_d(L_X\N\mathfrak d_1\N\mathfrak d_2/X)\), outside the
inner direct sum.  The part with this factor has direct scale
at least \(L_X/2\), and there are \(O(U^2)=O(X)\) pairs.
Its absolute value is therefore at most
\[
 C X\exp\left\{-c\frac{L_X^{1/4}}{\{\log(3L_X)\}^r}\right\}.
\]
Put \(t=\log(3X)\) and \(h=\log(e+t)\).
The exponent is at least \(c' t h^{\kappa/4-r}\) for large
\(X\).  Since \(\kappa>4r\), its ratio to \(\log X\) tends
to infinity.  This part is \(O_M(X^{-M})\) for every \(M\).
Only \(|\chi_d|\le1\) entered the direct estimate: no moving
inner weight was subjected to \eqref{jct:weight}.

The remaining expression is exactly \eqref{jct:rsoft}.
Its support gives
\(\N\mathfrak d_1\N\mathfrak d_2>X/L_X\),
\(\N\mathfrak c<2L_X\), and \(\N\mathfrak d_i<2U\).
Dividing by the other inverse norm yields
\(\N\mathfrak d_i>X/(2UL_X)=U/(4L_X)\).
This proves \eqref{jct:soft-main} and \eqref{jct:balanced-support}.

Fix \(\mathfrak c\), put \(n=\N\mathfrak c\),
\(Y_c=X/n\), and
\[
 V_c(t)=W(t)[1-\chi_d((L_X/n)t)].
\]
These functions have common fixed compact support and uniformly
bounded finite logarithmic-derivative seminorms: each derivative
acts on the argument of the changing cutoff.  Nearly analytic
bounds on this moving family are unnecessary.
The inner difference between \eqref{jct:rsoft} and
\eqref{jct:rjoint} is precisely
\(J_{V_c}(Y_c,U,Y_c^\gamma)\).
Here \(U^2/Y_c=2n\), and \(Y_c\ge X/(2L_X)\ge2\) for large
\(X\). By \eqref{jct:exponents},
\(\beta-a_\beta\gamma_{\beta,u}=1/2\).
Thus for every \(\theta>0\), Theorem~\ref{jct:tail-theorem} gives
\[
 |J_{V_c}(Y_c,U,Y_c^\gamma)|
 \ll_\theta Y_c^{1/2+\theta}\log(2+2n).
\]
Partial summation in \eqref{jct:divisor-count} gives
\[
 \sum_{\N\mathfrak c<2L_X}
 \frac{d_{2,K}(\mathfrak c)\log\N\mathfrak c
             \log(2+2\N\mathfrak c)}{(\N\mathfrak c)^{1/2}}
 \ll L_X^{1/2}\{\log(2L_X)\}^3.
\]
Sum over the outer ideal, choose \(\theta\) smaller than the
prescribed final \(\varepsilon\), and absorb these fixed powers
of logarithms.  This proves \eqref{jct:joint-main}; a bounded
initial range of \(X\) is absorbed in the constant.
For a surviving summand,
\(\N\mathfrak d_1\N\mathfrak d_2>Y_c\) and
\(r_1r_2<2Y_c^\gamma\), whence
\(q_1q_2>Y_c^{1-2\gamma}/4\).
This proves \eqref{jct:primitive-threshold}.

\section{The all-weight RH criterion}
At a split rational prime \(p\ne3,269\), the two values of
\(\psi\) over \(p\) sum to \(t(p)\).
Indeed, the cubic Kummer Frobenius is trivial precisely when
269 is a cube modulo \(p\); the two values then sum to two,
and otherwise sum to \(-1\).
At an inert prime the residue degree is two, so there is no
degree-one term.  Hence \(\Theta_W\) is exactly the
degree-one, first-power part of \(T_W\).
Every remaining rational norm is a prime power \(p^v\) with
\(v\ge2\).  There are at most two prime ideals over \(p\),
their logarithmic weights are at most \(2\log p\), and
there are \(O(\log(3X))\) possible exponents.  The elementary
bound \(p\le(2X)^{1/v}\) consequently gives
\begin{equation}
 T_W(X)-\Theta_W(X)
 \ll_W X^{1/2}\{\log(3X)\}^2.
 \label{jct:prime-powers}
\end{equation}
No prime number theorem is used here.

Under cubic Hecke RH, \(P\) is nonvanishing in \(\Re s>1/2\):
the finite factor \(H\) has zeros only on \(\Re s=0\).
Lemma~\ref{jct:reciprocal-growth}, with an Euler-product lower
bound on the line two, gives polynomial growth for \(1/P\)
on every bounded closed strip to the right of one half.
Cauchy's derivative formula on fixed-radius disks and the standard
strip growth of \(P\) give polynomial growth for \(P'\) there.
Their product \(-P'/P\) is therefore polynomially bounded.
Mellin inversion and rapid decay of \(\widehat W\) permit
moving its contour from two to \(1/2+\delta\), with no
residues and vanishing horizontal integrals.  Thus
\(T_W(X)\ll_{W,\delta}X^{1/2+\delta}\).
Equations \eqref{jct:joint-main} and \eqref{jct:prime-powers}
give the two stated square-root bounds.

Conversely, either asserted family of bounds implies
\(T_W(X)=O_{W,\varepsilon}(X^{1/2+\varepsilon})\) for every
indicated nonnegative weight and every \(\varepsilon>0\).
For \(\Re s>1\), direct integration gives the exact identity
\begin{equation}
 \int_0^\infty T_W(X)X^{-s-1}\,dX
 =\widehat W(s)\left(-\frac{P'}P(s)\right).
 \label{jct:converse-mellin}
\end{equation}
The integral near zero vanishes, since the shell contains no
nonzero von Mangoldt coefficient for sufficiently small \(X\).
The assumed bounds continue its left side holomorphically to
\(\Re s>1/2\), by dominated convergence on compact subsets,
choosing \(\varepsilon\) smaller than their distance from the
boundary.

Suppose \(L\) had a zero \(\rho=\beta+i\tau\) with
\(\beta>1/2\), of multiplicity \(m\ge1\).
Choose a sufficiently short open interval \(J\subset(1,2)\)
on which the phase \(\tau\log t\) varies by less than
\(\pi/3\).  The construction following
\eqref{jct:direct-decay} supplies a nonzero nonnegative admissible
weight supported in \(J\).  After one fixed rotation,
\(W(t)t^{\rho-1}\) has positive real part wherever \(W>0\),
so \(\widehat W(\rho)\ne0\).
Since \(H(\rho)\ne0\), the right side of
\eqref{jct:converse-mellin} then has a pole of residue
\(-m\widehat W(\rho)\), a contradiction.
Thus there are no nontrivial zeros to the right of one half;
the functional equation and conjugation symmetry give RH.
Central zeros lie on the boundary throughout, and every
multiplicity is included in the residue.
This proves Theorem~\ref{jct:rh-theorem} without discarding
zeros off the critical line at any intermediate step.

\section{Nonempty cancellation and its precise limitation}
At inverse norm of order \(X^{1/2}\), each individual rational
content is at most order \(X^{1/4}\), up to the fixed constants
and logarithmic ranges.  A one-variable content tail followed
by absolute summation in the other variable costs
\(X^{1+\varepsilon}R^{-a}\), with \(a=3/2\) or \(7/4\).
Making that bound square-root size requires a cutoff exponent
\(1/(2a)=1/3\) or \(2/7\), both greater than \(1/4\).
That particular individual-content truncation is empty at its
useful threshold.  The product of the two contents can instead
reach order \(Y_c^{1/2}\), so the joint cutoffs lie inside
its available range.

For an explicit nonzero removed summand, take \(u_0\in(1,2)\)
with \(W(u_0)>0\), and let \(\ell\) run through rational primes
other than 3 and 269.  Put
\[
 \mathfrak c=(2),\qquad \mathfrak d_1=\mathfrak d_2=(\ell),
 \qquad X=4\ell^4/u_0.
\]
The prime 2 is inert, so \(\N\mathfrak c=4\),
\(d_{2,K}(\mathfrak c)=2\), and \(\psi(\mathfrak c)=1\).
Both inverse norms equal \(\ell^2\), whereas
\(U=\sqrt{8/u_0}\,\ell^2>2\ell^2\); hence both caps equal
one.  The direct-cutoff argument is \(L_Xu_0/4\), whose
complementary factor is one for all sufficiently large \(\ell\).
The shell argument is \(u_0\), and
\(m((\ell))^2=\eta(\ell)^2=1\).
This is a nonzero summand of \(R_{\mathrm{soft},W}\).
Its joint content is \(\ell^2\), and
\[
 \frac{\ell^2}{(X/4)^\gamma}
 =u_0^\gamma\ell^{2-4\gamma}\longrightarrow\infty.
\]
Since \(0<\gamma_{\beta,u}<1/2\), it is removed by every stated
cutoff, while both
primitive norms are one.  This also verifies concretely why
cross-coprimality must not be added.

\paragraph{Strictness of the stronger cutoff.}
Let \(0<\gamma_1<\gamma_0\le1/3\) be any two of the cutoffs
compared in the statements, and choose
\(\gamma_1<\delta<\gamma_0\). Then
\(t=\delta/(1-2\delta)\) lies in \((0,1)\).
There are arbitrarily large split primes \(q\equiv1\pmod3\).
An elementary proof of their infinitude is to assume that all of
them divide an integer \(n\) which is also divisible by three:
a prime divisor of \(n^2+n+1\) is then new, is not three, and
has \(n\) of order three modulo it.
Choose a prime ideal \(\mathfrak q\) of norm \(q\).

For every sufficiently large \(R\), there is a squarefree integer
\(r\in[R,2R]\) prime to \(3\cdot269\). Indeed,
\(\mu^2(r)=\sum_{d^2\mid r}\mu(d)\) and elementary counting in
fixed residue classes give
\[
 \sum_{\substack{R<r\le2R\\(r,807)=1}}\mu^2(r)
 =c_{807}R+O(\sqrt R),\qquad
 c_{807}=\frac{\varphi(807)}{807}
         \prod_{p\nmid807}(1-p^{-2})>0.
\]
To check the error, sum the \(O(1)\) residue-count error over
\(d\le\sqrt{2R}\), and bound the remaining absolutely convergent
\(d^{-2}\) tail by \(O(R^{-1/2})\).
Take such an \(r\) with \(R=q^t\). Since \(t<1\), eventually
\(r<q\), and hence \((r,q)=1\).
Set
\[
 \mathfrak c=(2),\qquad
 \mathfrak d_1=\mathfrak d_2=(r)\mathfrak q,
 \qquad X=4(r^2q)^2/u_0.
\]
The two ideal coefficients are nonzero by
\eqref{jct:content-factorization}. Both caps equal one, the direct
cutoff equals its complementary value one for large \(q\), and
the shell argument equals \(u_0\). Thus this is a nonzero summand
of the balanced sum. Its content product is \(r^2\), and with
\(Y_c=X/4\),
\[
 \frac{\log(r^2)}{\log Y_c}\longrightarrow
 \frac{t}{2t+1}=\delta.
\]
Consequently \(r^2/Y_c^{\gamma_0}\to0\), whereas
\(r^2/Y_c^{\gamma_1}\to\infty\). The older cutoff is identically
one and the stronger cutoff identically zero on these summands
for all sufficiently large \(q\). This proves the literal strict
support assertion, rather than inferring it from exponents alone.
Moreover, every summand with a nonzero stronger cutoff eventually
has the older cutoff equal to one: its content product is less than
\(2Y_c^{\gamma_1}<Y_c^{\gamma_0}\), since
\(Y_c\ge X/(2L_X)\to\infty\). Thus the inclusion of supports
is asymptotically strict without any monotonicity assumption on
\(\chi\).

\paragraph{The numerical specialization.}
For \(\beta=B_*\), the formulas are
\[
 \gamma_{B_*,u}=\frac{B_*-1/2}{2B_*-u},\qquad
 1-2\gamma_{B_*,u}=\frac{1-u}{2B_*-u}.
\]
For fixed \(u\), the first rational function is strictly increasing in \(B_*\),
with derivative \((1-u)/(2B_*-u)^2>0\).
The exact algebraic description in \cite[equations~(42), (44)]{JCTLiu}
gives
\[
 \frac{8733458561067185768}{10^{19}}<B_*<
 \frac{8733458561067185769}{10^{19}}.
\]
For the middle row, in which \(u=B_*/2\) varies with \(B_*\),
the formulas simplify to
\[
 a_{B_*}=3B_*/2,\qquad
 \gamma_{B_*,B_*/2}=\frac{2B_*-1}{3B_*},\qquad
 1-2\gamma_{B_*,B_*/2}=\frac{2-B_*}{3B_*}.
\]
The cutoff has derivative \(1/(3B_*^2)>0\), while the last
expression is decreasing. Substitution of the rational endpoints
therefore proves all displayed approximations by rational arithmetic.
Reducing \(u\) from \(1/2\) to \(B_*/2\) increases the positive
denominator of the cutoff, so strictly decreases it.
For \(u=1/4\), clearing
the positive denominator shows
\[
 \gamma_{B_*,1/4}<\tfrac14
 \quad\Longleftrightarrow\quad B_*<\tfrac78.
\]
Similarly,
\[
 \frac{2B_*-1}{3B_*}<\frac27
 \quad\Longleftrightarrow\quad B_*<\frac78.
\]
For every fixed \(u\) in the table, monotonicity and \(B_*<1\) prove
\(\gamma_{B_*,u}<\gamma_{1,u}\). Also \(u<1\) gives
\(0<\gamma_{\beta,u}<1/2\) for every permitted \(\beta\).
These calculations prove Theorem~\ref{jct:liu-specialization}.

The secondary Euler factor also matters for the general conditional
theorem. Without extracting it the first correction product would
require a content real part exceeding \(1-\beta\), giving the
larger cutoff
\((\beta-1/2)/(2\beta-\max(u,1-\beta))\).
The two expressions agree when \(\beta+u\ge1\), which includes
every numerical row above. They differ strictly when
\(\beta+u<1\). For example, the conditional parameters
\(\beta=3/5,u=1/4\) give \(2/19\) in place of \(1/8\).
No zero-free theorem at \(\beta=3/5\) is asserted here.

For comparison, taking absolute values everywhere would give
only
\begin{equation}
 |J_V(Y,U,R)|\ll_{V,b,\chi}
 YR^{-1}\{\log(3Y)\}^2.
 \label{jct:absolute-comparison}
\end{equation}
For each product \(r=r_1r_2\), there are \(d(r)\) ordered
factorizations, and at most
\(O((Y/r^2)\log(2Y))\) primitive ideal pairs by counting.
Summing \(\sum_{r>R}d(r)/r^2\ll\log(2R)/R\), with
\(r<\sqrt{2Y}\), proves the bound.
At \(R=Y^{1/3}\), this is \(Y^{2/3}\) up to logarithms;
Theorem~\ref{jct:tail-theorem} gives square-root size up to
\(Y^\varepsilon\) and the displayed cap-aspect logarithm.
This is signed cancellation in a nonempty portion of the
balanced expression.  It does not bound the remaining portion.

\paragraph{Relation to other cancellation methods.}
The classical identity and elementary Mellin continuation are
the mechanisms proved here; no priority assertion is intended.
For context, the M\"obius estimate for fixed residue classes and
sectors in \cite[Lemma~5.2]{JCTLam} gives arbitrarily large
logarithmic savings for a fixed positive definite quadratic
form.  It does not supply the required square-root bound for
the remaining paired sum.
The estimate of \cite[Theorem~1.1]{JCTFG} concerns M\"obius
times cusp-form coefficients on the full modular group and
has exponential-logarithmic saving, not a fixed power saving.
The varying-polynomial result
\cite[Theorem~1.3]{JCTGMPP} averages over \(h\) in
\(n^2+h\); its theorem does not estimate the fixed balanced coefficient sum
here, even though differencing can create many dependent shifts.
None of these theorems is imported into this proof.

There is also a direct algebraic obstruction to treating the
fixed cubic character as generic bilinear oscillation.
For its finite-field realization \(\vartheta\) on
\(\mathbf F_{269^2}^{\times}\), the kernel
\(\vartheta(xy)=\vartheta(x)\vartheta(y)\) has rank one.
For nonzero \(y_1,y_2\), extending the character by zero at
zero gives exactly
\[
 \sum_{x\in\mathbf F_{269^2}}
 \vartheta(xy_1)\overline{\vartheta(xy_2)}
 =(269^2-1)\vartheta(y_1/y_2).
\]
Thus a generic Cauchy--Schwarz argument in the product variable
does not produce off-diagonal character cancellation.
Obtaining a square-root estimate for the remaining M\"obius
expression remains the substantive problem, precisely as
Theorem~\ref{jct:rh-theorem} records.

\section*{Method, record, and dedication}
The convolution algebra is attributed to Heath--Brown.  The content
factorization, contour estimates, and their application are proved
in this paper from the stated standard Hecke background.
The mathematics and text were produced by GPT-6 Astra in ChatGPT Pro
mode. The elementary joint tail, balanced identity and zero-detection
argument have complete separate statement-only reconstructions.
The stronger joint-tail estimate was independently re-proved from
its statement by another instance of the same model. The secondary
mixed Euler factor and the full conditional range have a further
complete statement-only reconstruction; the shorter factorization
from that reconstruction is used in the proof here. Its elementary
harmonic-function argument also supplies the reciprocal growth
bounds, without a separate zero-count input. The version
using two zero-free inputs was also derived separately before that
instance saw the proof here. A further same-model targeted check
examined the complete factorization, contour argument and strict
support comparison. The reconstructions explicitly grant the named
external analytic inputs and identify earlier proved dependencies.
The full arguments are printed here, without depending on another
unpublished note.
The effort setting was inherited and was not exposed in the review
record.  No human has checked this work.  Reproducible sources and
review records belong to the
\href{\workbenchurl}{workbench repository} and the
\href{\recordurl}{research record}.
The newly authored text and proof organization are dedicated to the
public domain under CC0 1.0.  Cited third-party works retain their own terms.
\begin{thebibliography}{9}
\bibitem[Heath--Brown(1982)]{JCTHeathBrown}
Heath--Brown, D. R. (1982).
Prime numbers in short intervals and a generalized Vaughan identity.
\emph{Canadian Journal of Mathematics, 34}(6), 1365--1377.
\url{https://doi.org/10.4153/CJM-1982-095-9}.
Printed p.~1367, Lemma~1, equations~(6)--(7).

\bibitem[Zaman(2016)]{JCTZaman}
Zaman, A. (2016).
Explicit estimates for the zeros of Hecke \(L\)-functions.
\emph{Journal of Number Theory, 162}, 312--375.
\url{https://doi.org/10.1016/j.jnt.2015.10.003}.
Primary preprint: \url{https://arxiv.org/pdf/1502.05679},
version~4, \S2.1, equations~(2.1)--(2.3) and
Lemmas~2.1--2.3, preprint pp.~6--7.

\bibitem[Liu(2026)]{JCTLiu}
Liu, B. (2026).
\emph{Joint marked moments and optimized zero-free half-planes for
unitary Hecke L-functions} (Version~2). Hexagon, 2611.00421v2.
\url{https://hexagonmath.org/2611.00421v2}.
Theorem~1.2, p.~3; equations~(42), (44), pp.~27--29;
Theorem~8.29, pp.~78--80. The analytic theorem is imported as a
stated external input; its 101-page proof is not independently
reconstructed here.

\bibitem[NIST(n.d.)]{JCTDLMF}
National Institute of Standards and Technology. (n.d.).
\emph{Digital Library of Mathematical Functions},
\S\S5.5 and 5.11.
\url{https://dlmf.nist.gov/5.5};
\url{https://dlmf.nist.gov/5.11}.
Formulas 5.5.1, 5.5.3, and 5.11.2.

\bibitem[Lam et al.(2020)]{JCTLam}
Lam, P. C.-H., Schindler, D., \& Xiao, S. Y. (2020).
On prime values of binary quadratic forms with a thin variable.
\emph{Journal of the London Mathematical Society, 102}(2), 749--772.
\url{https://doi.org/10.1112/jlms.12336}.
The comparison in this paper concerns Lemma~5.2.

\bibitem[Fouvry and Ganguly(2014)]{JCTFG}
Fouvry, E., \& Ganguly, S. (2014).
Strong orthogonality between the M\"obius function, additive
characters and Fourier coefficients of cusp forms.
\emph{Compositio Mathematica, 150}(5), 763--797.
\url{https://doi.org/10.1112/S0010437X13007732}.
Primary preprint: \url{https://arxiv.org/pdf/1301.5507}, Theorem~1.1.

\bibitem[Grimmelt et al.(2026)]{JCTGMPP}
Grimmelt, L., Merikoski, J., Pandey, M., \& Pascadi, A. (2026).
\emph{Primes and the M\"obius function over quadratic polynomials}
[Preprint]. arXiv.
\url{https://arxiv.org/abs/2610.41356}.
Version~1, Theorem~1.3 and \S1.1.
\end{thebibliography}
\end{document}
